\documentclass[10pt]{article}
\usepackage[a4paper,margin=2.1cm]{geometry}
\usepackage{fontspec}
\usepackage{amsmath,amssymb,booktabs,graphicx,hyperref,xcolor,caption,float,multirow,adjustbox,array}
\usepackage[section]{placeins}
\usepackage[numbers,sort&compress]{natbib}
\usepackage{xeCJK}
\IfFileExists{C:/Windows/Fonts/simsun.ttc}{
  \setCJKmainfont[AutoFakeBold=2]{SimSun}
  \setCJKsansfont{Microsoft YaHei}
  \setCJKmonofont{FangSong}
}{
  \setCJKmainfont{Noto Serif CJK SC}
  \setCJKsansfont{Noto Sans CJK SC}
  \setCJKmonofont{Noto Sans Mono CJK SC}
}
\renewcommand{\abstractname}{摘要}
\renewcommand{\figurename}{图}
\renewcommand{\tablename}{表}
\renewcommand{\refname}{参考文献}
\renewcommand{\contentsname}{目录}
\linespread{1.25}
\hypersetup{colorlinks=true,linkcolor=blue!45!black,citecolor=blue!45!black,urlcolor=blue!45!black}
\captionsetup{font=small,labelfont=bf,labelsep=quad}
\setlength{\emergencystretch}{3em}
\renewcommand{\arraystretch}{1.1}
\renewcommand{\topfraction}{0.92}\renewcommand{\bottomfraction}{0.85}\renewcommand{\textfraction}{0.08}\renewcommand{\floatpagefraction}{0.75}
\setcounter{topnumber}{4}\setcounter{totalnumber}{6}
\graphicspath{{../figures/}}
\newcommand{\NRows}{198}
\newcommand{\NFitProfiles}{9}
\newcommand{\NFitPoints}{533}
\newcommand{\NCheckPass}{36}
\newcommand{\NCheckAll}{36}
\newcommand{\ReproConsistent}{7}
\newcommand{\ReproAll}{12}
\newcommand{\SolverErr}{2\times10^{-5}}
\newcommand{\SolverCachedMs}{0.67}
\newcommand{\SolverLsodaMs}{33.3}
\newcommand{\SolverRadauMs}{1163}
\newcommand{\LandMinFrac}{46}
\newcommand{\LandMaxFrac}{94}
\newcommand{\TransferWins}{4}
\newcommand{\TransferN}{8}
\newcommand{\TransferRatio}{0.975}
\newcommand{\LotoWinsOrig}{5}
\newcommand{\LotoWinsFull}{4}
\newcommand{\LotoSigWinsOrig}{4}
\newcommand{\LotoSigLossOrig}{0}
\newcommand{\LotoSigWinsFull}{3}
\newcommand{\LotoSigLossFull}{1}
\newcommand{\LotoAllNew}{2.121}
\newcommand{\LotoAllOrig}{2.533}
\newcommand{\LotoAllFull}{2.312}
\newcommand{\LotoAllRedOrig}{16.2}
\newcommand{\LotoAllRedFull}{8.2}
\newcommand{\LotoAllPOrig}{$<$0.001}
\newcommand{\LotoAllPFull}{$<$0.001}
\newcommand{\LotoAkSeventyNew}{1.886}
\newcommand{\LotoAkSeventyOrig}{1.828}
\newcommand{\LotoAkSeventyFull}{1.655}
\newcommand{\LotoAkHundredNew}{3.505}
\newcommand{\LotoAkHundredOrig}{3.742}
\newcommand{\LotoAkHundredFull}{3.293}
\newcommand{\LotoBornNew}{1.054}
\newcommand{\LotoBornOrig}{1.825}
\newcommand{\LotoBornFull}{1.825}
\newcommand{\LotoTaoANew}{1.358}
\newcommand{\LotoTaoAOrig}{2.101}
\newcommand{\LotoTaoAFull}{1.440}
\newcommand{\LotoTaoBNew}{1.495}
\newcommand{\LotoTaoBOrig}{1.643}
\newcommand{\LotoTaoBFull}{1.591}
\newcommand{\LotoMaoNew}{3.098}
\newcommand{\LotoMaoOrig}{3.607}
\newcommand{\LotoMaoFull}{3.314}
\newcommand{\LotoBBGP}{5.67}
\newcommand{\LotoBBMLP}{12.99}
\newcommand{\LotoBBTabICL}{6.13}
\newcommand{\LotoBBNODE}{1.79}
\newcommand{\FcstWinsOrig}{4}
\newcommand{\FcstWinsFull}{4}
\newcommand{\FcstSigWinsOrig}{0}
\newcommand{\FcstSigLossOrig}{0}
\newcommand{\FcstSigWinsFull}{1}
\newcommand{\FcstSigLossFull}{1}
\newcommand{\FcstAllNew}{11.874}
\newcommand{\FcstAllOrig}{12.080}
\newcommand{\FcstAllFull}{12.060}
\newcommand{\FcstAllRedOrig}{1.7}
\newcommand{\FcstAllRedFull}{1.5}
\newcommand{\FcstAllPOrig}{0.184}
\newcommand{\FcstAllPFull}{0.740}
\newcommand{\FcstAkSeventyNew}{1.443}
\newcommand{\FcstAkSeventyOrig}{1.097}
\newcommand{\FcstAkSeventyFull}{9.527}
\newcommand{\FcstAkHundredNew}{5.506}
\newcommand{\FcstAkHundredOrig}{10.670}
\newcommand{\FcstAkHundredFull}{0.673}
\newcommand{\FcstBornNew}{3.971}
\newcommand{\FcstBornOrig}{4.044}
\newcommand{\FcstBornFull}{4.044}
\newcommand{\FcstTaoANew}{17.525}
\newcommand{\FcstTaoAOrig}{19.861}
\newcommand{\FcstTaoAFull}{19.861}
\newcommand{\FcstTaoBNew}{15.753}
\newcommand{\FcstTaoBOrig}{15.866}
\newcommand{\FcstTaoBFull}{15.866}
\newcommand{\FcstMaoNew}{7.765}
\newcommand{\FcstMaoOrig}{4.654}
\newcommand{\FcstMaoFull}{4.159}
\newcommand{\FcstBBGP}{28.05}
\newcommand{\FcstBBMLP}{40.79}
\newcommand{\FcstBBTabICL}{24.19}
\newcommand{\FcstBBNODE}{20.92}
\newcommand{\InAkSeventyFnew}{56.8}
\newcommand{\InAkSeventyFpaper}{59.0}
\newcommand{\InAkSeventyRed}{3.7}
\newcommand{\InAkHundredFnew}{31.8}
\newcommand{\InAkHundredFpaper}{117.9}
\newcommand{\InAkHundredRed}{73.0}
\newcommand{\InBornFnew}{105.0}
\newcommand{\InBornFpaper}{329.9}
\newcommand{\InBornRed}{68.2}
\newcommand{\InTaoAFnew}{75.1}
\newcommand{\InTaoAFpaper}{121.2}
\newcommand{\InTaoARed}{38.0}
\newcommand{\InTaoBFnew}{260.2}
\newcommand{\InTaoBFpaper}{325.8}
\newcommand{\InTaoBRed}{20.1}
\newcommand{\InMaoFnew}{1169.6}
\newcommand{\InMaoFpaper}{1475.9}
\newcommand{\InMaoRed}{20.8}
\newcommand{\InFobjWins}{6}
\newcommand{\InBICWins}{5}
\newcommand{\PoolAllSel}{2.117}
\newcommand{\PoolAllLo}{2.121}
\newcommand{\PoolAkHundredLii}{2.203}
\newcommand{\SimSMBiasOrig}{+0.57}
\newcommand{\SimSMBiasNew}{+0.05}
\newcommand{\SimSMBiasOrigPct}{76}
\newcommand{\SimSiiOrig}{1.213}
\newcommand{\SimSiiNew}{0.260}
\newcommand{\SimSiiiOrig}{1.739}
\newcommand{\SimSiiiFull}{0.943}
\newcommand{\SimSiiiNew}{0.718}
\newcommand{\SimSiOrig}{0.638}
\newcommand{\SimSiNew}{0.667}
\newcommand{\SimSiFull}{0.676}
\newcommand{\SimSiJoint}{0.794}
\newcommand{\NSimReps}{60}
\newcommand{\BornPM}{0.18}
\newcommand{\BornPD}{0.013}
\newcommand{\BornSM}{0.56}
\newcommand{\BornSD}{0.039}
\newcommand{\AkHundredKd}{0.018}
\newcommand{\TaoAKd}{0.001}
\newcommand{\AkSeventyKd}{6.3e-04}
\newcommand{\TaoAPD}{0.06}
\newcommand{\AkHundredFullGlobal}{48.9}
\newcommand{\MaoFullAtPaper}{1476.5}
\newcommand{\NLotoFolds}{145}
\newcommand{\AblLamOne}{2.209}
\newcommand{\AblJoint}{2.177}
\newcommand{\AblSimplOnly}{2.271}
\newcommand{\AblFullOnly}{2.109}
\newcommand{\AblNoProc}{2.289}
\newcommand{\AblNoProcBorn}{1.718}
\newcommand{\AblNoProcP}{$<$0.001}
\newcommand{\AblNoKd}{2.115}
\newcommand{\AblNoKdBorn}{1.052}
\newcommand{\AblNoKdP}{$<$0.001}
\newcommand{\AblNoBoth}{2.277}
\newcommand{\AblNoBothBorn}{1.716}
\newcommand{\AblNoBothP}{$<$0.001}
\newcommand{\HybLotoAll}{1.290}
\newcommand{\HybLotoBase}{2.121}
\newcommand{\HybLotoP}{$<$0.001}
\newcommand{\HybLotoMao}{1.197}
\newcommand{\HybLotoOrigGP}{1.307}
\newcommand{\HybFcstAll}{11.857}
\newcommand{\HybFcstBase}{11.874}
\newcommand{\HybLotoRedOrig}{49.1}
\newcommand{\HybLotoPOrig}{$<$0.001}
\newcommand{\HybLotoPOrigGP}{0.888}
\newcommand{\HybWinsOrig}{5}
\newcommand{\HybFcstOrig}{12.080}
\newcommand{\FoldN}{145}
\newcommand{\FoldWinsOrig}{108}
\newcommand{\FoldPOrig}{2.6\times10^{-10}}
\newcommand{\FoldWinsFull}{97}
\newcommand{\FoldPFull}{9.4\times10^{-6}}
\newcommand{\FoldHybWinsOrig}{126}

\newcommand{\inctab}[1]{\input{tables_zh/#1.tex}}
\newcommand{\incfig}[2][\linewidth]{\includegraphics[width=#1,height=0.62\textheight,keepaspectratio]{#2.pdf}}
\title{指纹信息先验与伪持续性动力学\\用于脂肪酶催化三羟甲基丙烷酯化的预测建模}
\author{稿件草稿；作者姓名与单位待补充}
\date{2026 年 10 月 2 日}
\begin{document}
\maketitle
\begin{abstract}脂肪酶催化三羟甲基丙烷（TMP）酯化涉及三个相互竞争的 Ping Pong bi bi 反应，有限的浓度观测难以同时约束底物选择性和动态速率参数。针对两步半机理方法固定指纹选择性、未显式描述伪持续性与批内失活的局限，本文提出带指纹信息先验的伪持续性动力学模型（PPK-FIP）。PPK-FIP 将指纹估计转化为可调宽度的先验，在联合参数估计中引入伪持续性与一阶失活，并通过精确时间重参数化降低模型求值成本。在四项研究的九条反应曲线上，留一时间点交叉验证的合并预测误差较重新拟合的初始简化模型降低 \LotoAllRedOrig\,\%，较初始全模型降低 \LotoAllRedFull\,\%（均为 $p<0.001$）。组件消融表明，预测改善主要来自 Bornadel 数据中的伪持续性修正，以及稀疏采样条件下的指纹先验约束。已知伪持续性真值的模拟中，初始两步法对选择性 $S_M$ 的高估为 \SimSMBiasOrigPct\,\%，PPK-FIP 将对数偏差降至 \SimSMBiasNew。高斯过程残差校正进一步将合并交叉验证误差降低至初始简化模型的约一半，但向前外推中 PPK-FIP 与初始简化模型的合并误差相近。研究结果支持在 TMP 酯化建模中联合评估组成路径、参数约束与样本外预测，并将伪持续性参数解释限定于经过验证的反应条件。
\end{abstract}
\noindent\textbf{关键词}\quad Ping Pong bi bi；脂肪酶；三羟甲基丙烷；生物润滑剂；动力学模型；选择性；伪持续性；酶失活；贝叶斯先验；交叉验证；混合建模
\section{引言}
脂肪酶催化酯化为生物质来源脂肪酸转化为润滑剂提供了反应路线。三羟甲基丙烷（TMP）经逐级酯化生成单酯（M）、二酯（D）和三酯（T），产物组成影响润滑剂的性能，因而反应器设计需要同时描述转化速率与酯类分布 \cite{akerman2011,mao2024}。近年来，Montiel 等研究了无溶剂条件下支链酸与支链醇的酶催化生物润滑剂合成及催化剂重复使用 \cite{montiel2024}；Mao 等将油酸裂解所得壬酸用于 TMP 酯合成，并评价了脂肪酶批次稳定性和产物润滑性能 \cite{mao2024}。可再生资源制备 TMP 的研究进一步拓展了全生物基原料路线 \cite{sayed2025}。合成路线与催化剂性能的进展使动力学模型的预测能力成为过程放大与条件选择的重要依据。

固定化酶反应的表观动力学还受到载体传质、底物吸附与酶稳定性的共同影响。Sengupta 等将竞争吸附和多孔树脂表征结合，研究了固定化脂肪酶催化酯化的吸附与孔扩散条件 \cite{sengupta2026}。Ribeiro 等在壳聚糖--黄原胶双层纳米胶囊固定化脂肪酶中考察了热稳定性，并用一阶模型描述主要热失活过程 \cite{ribeiro2025}。固定化脂肪酶的吸附与热失活研究分别说明了传质条件和活性衰减在特定固定化酶体系中的作用，但不能直接确定 TMP 连续酯化中的伪持续性概率或失活常数。TMP 动力学分析需要利用组分轨迹区分组成路径偏差与反应时间尺度偏差。

三个连续酯化步骤可由共享分母的 Ping Pong bi bi 方程描述 \cite{cleland1963,cornishbowden2012,bornadel2013}。完整速率方程包含多个相互关联的参数，有限浓度--时间观测往往不足以约束全部参数 \cite{alhaque2012}。Nishimura 等提出了两步半机理估计方法，本文称为初始两步法 \cite{nishimura2026}。指纹步骤以反应度 $c$ 为自变量，估计第二和第三次酯化的选择性 $S_M$、$S_D$；动态步骤固定选择性，再估计集总速率参数。公共速率因子在反应度域中消去，是指纹方法分离组成路径与时间尺度的基础 \cite{mitchell2022a,mitchell2022b,mitchell2023}。Nishimura 等利用四项研究的公开曲线，以残差平方和和贝叶斯信息准则（BIC）比较了简化与完整动态模型。

初始两步法的选择性约束仍存在两个需要检验的问题。不可逆指纹模型未包含伪持续性，即新生成的酯在释放前再次酰化的过程。Nishimura 等已指出，Bornadel 数据的传质限制可能使指纹参数仅代表表观选择性 \cite{nishimura2026,mitchell2023}。固定存在模型偏差的指纹估计，会限制动态步骤修正组分轨迹的能力。初始动态模型采用批次独立的速率常数描述重复使用过程，但未显式估计批内失活。此外，Nishimura 等的模型比较主要基于样本内拟合，BIC 参数计数未包括第一步从相同数据估计的两个选择性。组成路径修正能否改善留出观测预测，以及参数增加是否造成过拟合，需要独立于全数据拟合的检验。

近期酶反应网络与科学机器学习研究为参数估计和验证提供了相关思路。van Sluijs 等将最优实验设计与流动化学结合，通过迭代构建训练数据提高酶反应网络的动力学信息量 \cite{vansluijs2024}。Giampiccolo 等将部分已知动力学嵌入混合神经常微分方程，同时考察机理参数估计与可辨识性 \cite{giampiccolo2024}。de Rooij 等使用生理约束正则化改善通用微分方程在稀疏数据下的训练与预测 \cite{derooij2025}。酶反应网络实验设计与混合动力学研究表明，数据覆盖范围、机理约束和参数辨识需要联合评估，较低训练误差本身不能替代预测验证。表格基础模型 TabPFN 和 TabICL 也扩展了小样本数据建模的方法选择 \cite{hollmann2025,qu2025tabicl}，但通用表格预测能力并不直接等同于受质量守恒约束的反应外推能力。

本文研究指纹选择性约束与机理缺项对 TMP 酯化预测的影响，并提出带指纹信息先验的伪持续性动力学模型（PPK-FIP）。PPK-FIP 保留反应度域的选择性估计，将固定选择性改为经过校准的先验，并在动力学方程中显式加入伪持续性和一阶失活。精确时间重参数化用于降低多起点估计与交叉验证的计算成本。表~\ref{tab:roadmap} 概括了初始模型评估、组分残差诊断、样本外预测、组件消融、参数恢复及混合建模的检验设计。所有方法在相同公开数据上比较，指纹估计在每个验证折内重新计算；跨条件迁移和已知真值模拟用于限定机理参数的解释范围。

\begin{table}[!htbp]
\centering\small
\caption{TMP 酯化动力学模型的六项检验及对应证据。图~\ref{fig:workflow}、\ref{fig:data} 和 \ref{fig:bench} 给出估计流程、数据与计算性能。}\label{tab:roadmap}
\begin{tabular}{@{}c p{5.1cm} p{4.4cm} p{4.4cm}@{}}
\toprule
步骤 & 检验目标 & 分析方法 & 对应图表 \\
\midrule
1 & 初始两步模型的拟合一致性、参数可辨识性与样本内评价局限 & 公开数据核对、两步参数重新估计、多起点搜索与剖面目标函数分析 & 表~\ref{tab:data}--\ref{tab:reprodyn}；图~\ref{fig:repro}、\ref{fig:landscape} \\
2 & 初始两步模型的组分浓度残差与伪持续性、批内失活的关系 & 全数据拟合、组成路径比较与组分残差诊断 & 表~\ref{tab:insample}、\ref{tab:params}；图~\ref{fig:insample}--\ref{fig:resid} \\
3 & PPK-FIP 对留出时间点及后期反应轨迹的预测能力 & 留一时间点交叉验证与前半段训练的向前外推 & 表~\ref{tab:loto}、\ref{tab:fcst}；图~\ref{fig:loto}--\ref{fig:fcst} \\
4 & 伪持续性、失活、指纹先验及模型平均对预测误差的贡献 & 组件消融、先验宽度敏感性与跨条件参数池化 & 表~\ref{tab:ablation}、\ref{tab:pooling}；图~\ref{fig:ablation}、\ref{fig:pooling} \\
5 & 伪持续性组成路径的跨条件适用性与选择性估计偏差 & 留出条件的组成路径预测与蒙特卡洛参数恢复 & 表~\ref{tab:transfer}、\ref{tab:sim}；图~\ref{fig:transfer}、\ref{fig:sim} \\
6 & 机理约束与高斯过程残差校正对内插和外推的作用 & 数据驱动基线对比与机理--高斯过程混合建模 & 表~\ref{tab:lotobb}--\ref{tab:hybrid}；图~\ref{fig:blackbox}、\ref{fig:hybrid} \\
\bottomrule
\end{tabular}
\end{table}

\section{材料与方法}\label{sec:methods}
\subsection{初始两步法与 PPK-FIP 的估计流程}
图~\ref{fig:workflow} 比较了 Nishimura 等的初始两步法 \cite{nishimura2026} 与 PPK-FIP。初始两步法先用不可逆指纹模型估计 $S_M$、$S_D$，再固定选择性，用局部搜索估计动态参数，并根据样本内残差平方和 $F_\mathrm{obj}$ 和 BIC 选择集总常数形式。PPK-FIP 在组成路径中加入伪持续性，在时间映射中加入批内失活，并通过多起点最大后验（MAP）估计联合优化参数。指纹估计作为选择性先验的中心，先验宽度允许时间域数据修正选择性；收缩先验限制伪持续性与失活参数的幅度。简化和完整集总常数形式的预测按 BIC 权重平均，超参数通过留一案例元验证选择。初始两步法与 PPK-FIP 使用相同的训练与验证划分。
\begin{figure}[!htbp]\centering\incfig{fig01_workflow}\caption{初始两步法 \cite{nishimura2026}（a）与 PPK-FIP（b）的估计流程。指纹估计在初始两步法中固定，在 PPK-FIP 中用作先验中心。底部给出共同的评价流程。}\label{fig:workflow}\end{figure}

\subsection{反应数据与一致性核对}
反应曲线由程序从 Nishimura 等 \cite{nishimura2026} 的补充表 S1--S6 和 S17--S21 提取。Åkerman 数据包括 70 和 100\,°C 下的油酸--TMP 酯化，催化剂为 Novozym 435，初始酸与 TMP 摩尔比为 3:1 \cite{akerman2011}。Bornadel 数据采用相同底物体系 \cite{bornadel2013}。Tao 数据采用辛酸和 \emph{Candida} sp.\ 99--125 脂肪酶，初始摩尔比为 3.2:1，水活度分别为 0.25、0.35 和 0.45 \cite{tao2015}。Mao 数据采用壬酸与相同来源的脂肪酶，动态拟合使用前三个重复使用批次 \cite{mao2024}。组分浓度以相对于初始 TMP 的摩尔百分数表示，反应度定义为 $c=M+2D+3T$。公开数据共包含 11 条曲线、\NRows{} 行；与 Nishimura 等相同的 \NFitProfiles{} 条拟合曲线包含 \NFitPoints{} 个组分观测（表~\ref{tab:data}）。
\inctab{data}

图~\ref{fig:data} 给出九条拟合曲线在时间域和反应度域中的分布。Åkerman 曲线仅含 7--9 个非零采样时间，Bornadel 曲线含 30 个，采样密度差异影响参数约束强度。Tao 的三个水活度及 Mao 的三个批次在时间域中具有不同反应速率，但在反应度域中呈现接近的组成路径，因而适合检验路径的跨条件迁移。Mao 数据未报告 Z，且第一批次的三酯观测在约 90\,\% 处形成平台。有限时间平台与完全转化的渐近状态需要区分，平台偏差的来源不能仅由公开曲线确定。
\begin{figure}[!htbp]\centering\incfig{fig02_data_overview}\caption{九条拟合曲线在时间域（每组左图）和反应度域（每组右图）中的分布。Z、M、D、T 分别表示 TMP、单酯、二酯和三酯。数据取自文献 \cite{nishimura2026} 的补充表 S1--S6、S17--S19，单位为相对初始 TMP 的摩尔百分数。}\label{fig:data}\end{figure}

Nishimura 等未直接列出各拟合的观测数，本文通过依赖样本量的统计量核对数据范围。12 个动态拟合的 BIC 用于检查观测数 $n_T$ 与参数计数，4 个指纹拟合的残差平方和用于检查公开数据与报告拟合的一致性。动态参数核对包括将文献参数代入方程求值，以及对相同模型进行 30 起点重新估计。200 个随机起点的局部解分布用于评估优化敏感性，沿 $F_A$ 的剖面目标函数通过重新优化其他参数得到，用于检查当前数据下的参数可辨识性 \cite{raue2009}。

\subsection{初始两步模型及比较基线}
在持续移除生成水的不可逆酯化假设下，Nishimura 等的指纹模型为 \cite{nishimura2026}
\begin{equation}
\begin{aligned}
&\frac{dZ}{dc}=\frac{-Z}{S},\qquad \frac{dM}{dc}=\frac{Z-S_MM}{S},\qquad \frac{dD}{dc}=\frac{S_MM-S_DD}{S},\qquad\frac{dT}{dc}=\frac{S_DD}{S},\\
&S=Z+S_MM+S_DD.
\end{aligned}\label{eq:fp}
\end{equation}
动态模型使用公共 Ping Pong 速率因子
\begin{equation}
\begin{aligned}
&r=\frac{k_A[E]_T\,A}{A(F_A+F_{AZ}Z+F_{AM}M+F_{AD}D)+Z+S_MM+S_DD},\\
&\frac{dZ}{dt}=-rZ,\quad \frac{dM}{dt}=r(Z-S_MM),\quad \frac{dD}{dt}=r(S_MM-S_DD),\quad \frac{dT}{dt}=rS_DD.
\end{aligned}\label{eq:dyn}
\end{equation}
其中 $A=A_0-c$ 表示脂肪酸。指纹步骤通过最小二乘估计 $S_M$、$S_D$；动态步骤固定选择性，估计 $k_E=k_A[E]_T$ 与集总常数。简化形式采用 Nishimura 等选择的非零常数组合，完整形式保留全部集总常数。由于 $k_A$ 与 $[E]_T$ 仅以乘积出现，本文直接估计 $k_E$，并用 $k_E/k_A$ 核对文献隐含的酶载量。

目标函数为各组分浓度的残差平方和 $F_\mathrm{obj}=\sum(y_\mathrm{pred}-y_\mathrm{obs})^2$，Mao 数据的求和不包括 Z。BIC 定义为 $n_T\ln(F_\mathrm{obj}/n_T)+k\ln n_T$ \cite{schwarz1978}。为与文献比较，BIC 沿用 Nishimura 等仅计算动态步骤参数的约定；另报告计入两个指纹选择性的 $\mathrm{BIC}_\mathrm{tot}$。ORIG 与 ORIG-full 分别表示初始简化模型和初始全模型在本文中重新估计的基线。多起点搜索的最小目标值用于后续比较；有限起点搜索不构成全局最优的数学保证。

\subsection{伪持续性与批内失活的机理扩展}
伪持续性描述新生成的酯在扩散离开活性位点之前再次酰化的有效过程 \cite{mitchell2023}。设 $p_M$ 为新生成单酯在释放前继续酯化的概率，$p_D$ 为新生成二酯继续酯化的概率。以公共速率因子 $r$ 归一化的组成变化为
\begin{align}
\dot Z&=-Z, & \dot M&=Z(1-p_M)-S_MM,\nonumber\\
\dot D&=Zp_M(1-p_D)+S_MM(1-p_D)-S_DD, & \dot T&=Zp_Mp_D+S_MMp_D+S_DD,\nonumber\\
\Sigma&\equiv\dot c=Z(1+p_M+p_Mp_D)+S_MM(1+p_D)+S_DD.\label{eq:proc}
\end{align}
当 $p_M=p_D=0$ 时，反应度域方程退化为式~(\ref{eq:fp})。批内失活通过一阶活性衰减 $r\to r\,e^{-k_dt}$ 描述。参数 $k_d=0$ 对应无失活极限。因此，机理扩展允许通过约束 $p_M$、$p_D$ 和 $k_d$ 构造嵌套对照模型。

\subsection{精确时间重参数化}
共享速率因子使组分方程可写为反应度 $c$ 的函数，即将式~(\ref{eq:proc}) 除以 $\Sigma$。时间映射由求积关系给出
\begin{equation}
k_E\,\tau(c)=J_0(c)+F_AJ_1(c)+F_{AZ}J_Z(c)+F_{AM}J_M(c)+F_{AD}J_D(c),\qquad \tau=\frac{1-e^{-k_dt}}{k_d}.\label{eq:tau}
\end{equation}
其中 $J_0=\int_0^c S/((A_0-c')\Sigma)\,dc'$，$J_1=\int_0^c dc'/\Sigma$，$J_X=\int_0^c X/\Sigma\,dc'$（$X=Z,M,D$）；$k_d=0$ 时取极限 $\tau=t$。当 $p_M=p_D=0$ 时，$J_0=\ln[A_0/(A_0-c)]$。组成方程与积分由 Numba 编译的四阶 Runge--Kutta 方法求解，积分网格在接近完全转化时加密 \cite{lam2015}。给定组成路径后，任意时刻的预测由式~(\ref{eq:tau}) 的单调插值获得；更新 $k_E$、$F_\ast$ 或 $k_d$ 无需重新积分组成路径。

图~\ref{fig:bench} 比较 Bornadel 采样设计下带失活的伪持续性模型的求值成本。LSODA 与高精度 Radau 的单次求值分别需要约 19\,ms 和 \SolverRadauMs\,ms；重参数化方法在重新计算组成路径时需要约 11\,ms，使用缓存路径时需要 \SolverCachedMs\,ms。相对于容差 $10^{-12}$ 的 Radau 参考解，最大组分误差为 $\SolverErr$ 摩尔\,\%。缓存路径降低了动态参数搜索成本，使 30 起点估计、\NLotoFolds{} 折交叉验证和每情景 \NSimReps{} 次模拟的计算规模可控。
\begin{figure}[!htbp]\centering\incfig[0.75\linewidth]{fig05_solver_benchmark}\caption{Bornadel 采样设计下的模型求值时间与数值误差。重参数化方法分别统计重新计算组成路径与缓存组成路径的成本。误差相对于容差 $10^{-12}$ 的 Radau 参考解计算。}\label{fig:bench}\end{figure}

\subsection{指纹信息先验与模型平均}
固定指纹选择性可视为选择性先验宽度趋于零的极限，无正则化联合拟合则不施加指纹约束。PPK-FIP 采用有限宽度的指纹信息先验
\begin{equation}
\ln S_X\sim\mathcal N\!\left(\ln\hat S_X^{\,\mathrm{fp}},\ (\lambda\,\delta_X/1.96)^2\right),\quad p_M,p_D\sim\mathcal N^+(0,s_p^2),\quad \log_{10}k_d\sim\mathcal N(-3,1.5^2).\label{eq:prior}
\end{equation}
其中 $\delta_X$ 为训练数据上指纹估计的相对 95\,\% 置信半宽，$\lambda$ 控制选择性先验宽度，$s_p$ 控制伪持续性收缩强度。MAP 估计将先验残差 $\hat\sigma(u-\mu)/s$ 附加到数据残差向量，$\hat\sigma$ 为无惩罚拟合的残差标准差。优化使用信赖域反射算法，起点来自无惩罚联合模型及加入失活、伪持续性或两种机理的嵌套模型。

简化和完整集总常数形式分别拟合，预测按 $w_i\propto e^{-\mathrm{BIC}_i/2}$ 加权平均 \cite{hoeting1999,burnham2002}。每个目标案例的超参数 $(\lambda,s_p)\in\{1,3\}\times\{0.05,0.2\}$ 由留一案例元验证选择，选择准则为排除目标案例后的五个案例上 PPK-FIP 与 ORIG 的交叉验证误差比均值。目标案例不参与超参数选择，但同一研究来源可能包含多个案例，留一案例不等同于留一独立研究。指纹先验由训练数据估计，属于数据自适应正则化；先验与似然并非来自相互独立的数据。

\subsection{预测验证与组件分析}
\textbf{样本内拟合。}各方法使用全部训练观测进行 30 起点估计，报告 $F_\mathrm{obj}$、BIC、$\mathrm{BIC}_\mathrm{tot}$ 与 AICc \cite{hurvich1989}，并比较残差及反应度域组成路径。

\textbf{留一时间点交叉验证。}对 \NLotoFolds{} 个非零采样时间逐一留出相应时刻的全部组分观测，在训练数据上重新估计各方法的全部参数，再预测留出观测。初始两步法的指纹步骤与动态步骤均在每个验证折内重新估计。

\textbf{向前外推。}各方法使用 $t\le0.5\,t_\mathrm{max}$ 的观测训练，预测后半段反应曲线。每个案例使用一次时间划分，用于比较给定训练窗口下的外推表现。

\textbf{组件消融与多任务池化。}分别去除伪持续性、失活、指纹先验、模型平均或超参数元选择，每个变体均重新进行完整交叉验证。多任务扩展比较共享组成路径（L1）、进一步对集总常数施加层次收缩（L1h）及完全共享集总常数（L2），池化层级由留一案例元验证选择。

\textbf{跨条件迁移。}使用同一反应体系的其他温度、水活度或批次曲线估计组成路径，并在反应度域预测留出条件的组分浓度。无伪持续性路径与伪持续性路径使用相同迁移划分。

\textbf{蒙特卡洛参数恢复。}基于三个真实案例的采样设计与残差噪声生成合成数据。S1 使用初始无伪持续性模型，S2 包含伪持续性，S3 包含批内失活，每种情景重复 \NSimReps{} 次。曲线误差在稠密时间网格上相对于无噪声真值计算，选择性偏差在对数尺度上计算。

\textbf{数据驱动基线与混合模型。}比较高斯过程（Matérn-5/2 核加白噪声）、五成员多层感知机（MLP）集成、预训练 TabICL 回归器和神经常微分方程（神经 ODE）\cite{rasmussen2006,qu2025tabicl,chen2018}。输入包括时间、条件描述量及已知初始状态。PPK-FIP-GP 在机理预测上叠加按组分估计的高斯过程训练残差模型，零均值偏差项在远离训练数据时衰减。ORIG+GP 使用相同残差校正结构作为对照。

\subsection{统计分析与计算环境}
预测误差差异采用配对平方误差的 Wilcoxon 符号秩检验，六个案例之间采用 Holm 校正 \cite{wilcoxon1945,holm1979}。同时报告以时间折为单位的比较，以减少同一时刻多组分观测的相关性对检验的影响。交叉验证折仍共享部分训练数据，显著性结果需结合效应量与案例差异解读。计算使用 Python、NumPy、SciPy、scikit-learn 和 PyTorch，既有实验在 NVIDIA DGX Spark（GB10）工作站完成 \cite{virtanen2020}。数值与表格由保存的结果文件生成。

\section{结果}
\subsection{初始两步模型的拟合一致性与参数可辨识性}\label{sec:repro}
公开数据的观测范围与 Nishimura 等报告的拟合统计量总体一致 \cite{nishimura2026}。表~\ref{tab:datacheck} 中六个案例的观测数分别为 36、28、120、64、144 和 141，将观测数与文献报告的 $F_\mathrm{obj}$、参数个数代入 BIC 公式后，12 个 BIC 均与报告精度一致。四个指纹拟合的残差平方和分别为 74.2、840.6、535.3 和 1616.0，全部 \NCheckPass{}/\NCheckAll{} 项核对通过。统计量一致为数据提取和实现提供了交叉检查，但统计量相同本身不能证明逐个观测值完全一致。
\inctab{datacheck}

指纹选择性估计与文献报告值接近（表~\ref{tab:reprofp}，图~\ref{fig:repro}a）。重新估计的 $S_M$ 与报告精度一致，$S_D$ 的差别主要出现在报告位数与舍入精度上；95\,\% 置信半宽的相对差异约为 2--12\,\%。置信区间差异可能与 Hessian 的数值近似有关。后续 PPK-FIP 使用相同指纹方程获得先验中心，并在每个交叉验证训练集上重新估计。
\inctab{repro_fp}

动态参数核对显示，12 个拟合中有 \ReproConsistent{} 个在文献参数处得到的目标值与报告值相差不超过 1.5\,\%（表~\ref{tab:reprodyn}）。Åkerman 与 Bornadel 案例的酶载量分别接近 2.0\,\% 和 5.0\,\%，Mao 案例接近 15.0\,\%（图~\ref{fig:repro}c）。Mao 全模型采用实际初始酸与 TMP 摩尔比 10:3.3（$A_0=303.0$）时，目标值为 \MaoFullAtPaper，与文献报告的 1475.9 接近。Tao 案例的四组文献参数以及 Mao 简化模型的文献参数未满足 1.5\,\% 的一致性阈值，说明公开参数与本文实现之间仍有未解释的差异。

多起点重新估计进一步降低了部分初始全模型的目标值（图~\ref{fig:repro}b）。Åkerman 100\,°C 案例的 $F_\mathrm{obj}$ 从文献报告的 117.9 降至 \AkHundredFullGlobal；Tao 与 Mao 全模型也获得更低目标值。Tao 简化模型重新估计的目标值仍高于文献报告值，因此不能将全部差异归结为局部优化。后续比较采用本文在统一方程、初始条件与观测集合上重新估计的 ORIG 和 ORIG-full。
\inctab{repro_dyn}
\begin{figure}[!htbp]\centering\incfig{fig03_reproduction}\caption{初始两步模型的数值一致性与优化结果。（a）文献报告值与重新计算值的比较。（b）文献报告的目标值、文献参数处的目标值及多起点搜索最小值；F 和 S 分别表示完整与简化模型。（c）由 $k_E/k_A$ 估计的酶载量。文献报告值来自 Nishimura 等 \cite{nishimura2026}。}\label{fig:repro}\end{figure}

初始全模型对优化起点敏感（图~\ref{fig:landscape}a）。200 个随机起点中，仅 \LandMinFrac--\LandMaxFrac\,\% 到达已找到的最低目标值，部分局部解的目标值高出数个数量级。图~\ref{fig:landscape}b 中三个案例的 $F_A$ 剖面在搜索范围内持续下降或趋于平坦，未形成封闭的最优区间。在大 $F_A$ 极限下，速率主要受 $k_E$ 与集总常数的比值控制。剖面结果表明，当前观测对完整参数集的实际辨识能力有限 \cite{raue2009,transtrum2015}；仅凭有限数据上的剖面曲线，尚不能认定模型在所有实验设计下均结构不可辨识。
\begin{figure}[!htbp]\centering\incfig{fig04_landscape_identifiability}\caption{初始全模型的优化敏感性与参数剖面。（a）200 个随机起点的局部最小值及文献报告值（菱形）；上方数字为达到已找到最低目标值的起点比例。（b）沿 $F_A$ 的剖面目标函数，其他参数重新优化。剖面未在搜索范围内形成封闭最优区间。}\label{fig:landscape}\end{figure}

样本内统计量同时受到优化结果、参数计数与可辨识性的影响。因此，本文以多起点重新估计的基线开展模型比较，同时报告计入选择性的 $\mathrm{BIC}_\mathrm{tot}$，并用留出观测预测评价机理扩展。

\subsection{两步模型的组分浓度失配与伪持续性、批内失活}
PPK-FIP 在六个案例中的全数据残差平方和均低于初始简化模型和初始全模型（表~\ref{tab:insample}，图~\ref{fig:insample}）。相对于 Nishimura 等各案例报告的最小目标值，PPK-FIP 的比值分别为 0.96、0.27、0.32、0.62、0.80 和 0.79。计入全部估计参数后，PPK-FIP 的 $\mathrm{BIC}_\mathrm{tot}$ 在六个案例中的 \InBICWins{} 个低于 ORIG，差值约为 14--123；Åkerman 70\,°C 的 $\mathrm{BIC}_\mathrm{tot}$ 增加约 6，表明小幅拟合改善不足以抵消参数惩罚。
\inctab{insample}
\begin{figure}[!htbp]\centering\incfig{fig08_insample}\caption{全数据拟合比较。（a）各模型 $F_\mathrm{obj}$ 与文献 \cite{nishimura2026} 报告最小值的比值。（b）相对于 ORIG 的 $\mathrm{BIC}_\mathrm{tot}$ 差值，负值表示更低的信息准则。}\label{fig:insample}\end{figure}

Åkerman 两条温度曲线的拟合改善幅度明显不同。70\,°C 下 PPK-FIP 的 $F_\mathrm{obj}$ 为 \InAkSeventyFnew，与初始全模型接近，伪持续性与失活参数未显示强约束。100\,°C 下的 $F_\mathrm{obj}$ 从文献报告的 117.9 降至 \InAkHundredFnew，差异主要出现在后期反应减慢阶段。失活常数的模型加权值为 $k_d=\AkHundredKd$\,h$^{-1}$，占优简化形式的估计约为 0.032\,h$^{-1}$（表~\ref{tab:params}）。一阶失活有助于描述高温曲线，但浓度轨迹不能替代独立酶活测量。

Bornadel 数据的改善集中在早期单酯与二酯分布。初始模型在前两小时高估 M 约 2--5 摩尔\,\%，低估 D 约 4--5 摩尔\,\%，随后残差符号发生改变（图~\ref{fig:resid}b）。加入伪持续性后，早期残差缩小至约 $\pm2$ 摩尔\,\%，$F_\mathrm{obj}$ 从 329.9 降至 \InBornFnew。估计的 $p_M=\BornPM$，线性化 95\,\% 置信半宽约为 0.025，$p_D=\BornPD$。反应度域的比较也显示，伪持续性路径改善了 M 与 D 的相对分布（图~\ref{fig:paths}）。$S_M$、$S_D$ 分别由初始指纹估计的 0.909、0.066 调整至 \BornSM、\BornSD，与持续性指纹分析的解释方向一致 \cite{mitchell2023}。由于公共速率因子在反应度域中消去，仅调整动态集总常数不能修正给定选择性下的组成路径。

Tao 数据的拟合改善涉及完整集总常数形式和较小的二酯伪持续性参数。$a_w=0.25$ 时，$F_\mathrm{obj}$ 从文献报告的 121.2 降至 \InTaoAFnew，简化形式的权重仅为 0.02，$p_D=\TaoAPD$。$a_w=0.35/0.45$ 时，目标值从 325.8 降至 \InTaoBFnew。表~\ref{tab:params} 中三个水活度的 $p_D$ 均约为 0.05--0.06，组成路径变化小于 Bornadel 数据中的变化。

Mao 三个批次使用批次独立的 $k_E$，PPK-FIP 的 $F_\mathrm{obj}$ 从文献报告的 1475.9 降至 \InMaoFnew。较大的拟合残差仍集中在第一批次后期的 T 与 D（图~\ref{fig:fits}）。70\,h 后的 T 观测约为 90\,\%，D 约为 5--10\,\%，与当前不可逆模型的后期轨迹存在偏差。公开曲线不足以区分测量或数字化误差、终点未达完全转化以及未建模反应因素，因而不能将平台直接判定为数据误差。图~\ref{fig:resid}a 显示，PPK-FIP 在六个案例中的绝对残差中位数均低于 ORIG，但 Mao 的较大残差未被完全消除。
\inctab{params}
\begin{figure}[!htbp]\centering\incfig{fig06_fits_time}\caption{九条反应曲线的全数据拟合。标记表示观测，虚线表示重新估计的初始两步法，实线表示 PPK-FIP。}\label{fig:fits}\end{figure}
\begin{figure}[!htbp]\centering\incfig{fig07_paths_c}\caption{反应度域中的组分浓度路径。虚线表示初始指纹路径，实线表示含伪持续性的 PPK-FIP 路径。面板标题给出伪持续性估计，颜色与图~\ref{fig:fits} 一致。}\label{fig:paths}\end{figure}
\begin{figure}[!htbp]\centering\incfig{fig09_residuals}\caption{组分浓度残差比较。（a）六个案例的全数据绝对残差。（b）Bornadel 数据中 M 与 D 的残差。圆点和虚线对应初始模型，方块和实线对应 PPK-FIP。}\label{fig:resid}\end{figure}

组分残差与参数估计支持两种有条件的机理解释。Bornadel 数据的 M--D 分布偏差与伪持续性相符，Åkerman 100\,°C 的后期速率偏差可由批内失活缓解。全数据拟合尚不能判断新增参数是否改善未参与估计的观测预测。

\subsection{PPK-FIP 的留出时间点预测改善与外推边界}
留一时间点交叉验证中，各方法在全部 \NLotoFolds{} 折内重新估计参数（表~\ref{tab:loto}，图~\ref{fig:loto}）。在 \NFitPoints{} 个留出组分观测上，PPK-FIP 的合并 RMSE 从 ORIG 的 \LotoAllOrig{} 降至 \LotoAllNew{} 摩尔\,\%，相对降幅为 \LotoAllRedOrig\,\%（$p$\,\LotoAllPOrig）；相对于 ORIG-full 的 \LotoAllFull，降幅为 \LotoAllRedFull\,\%（$p$\,\LotoAllPFull）。PPK-FIP 的 RMSE 在 \LotoWinsOrig{} 个案例中低于 ORIG，其中 \LotoSigWinsOrig{} 个达到 Holm 校正后的显著性阈值。Bornadel、Tao 与 Mao 的改善较为明显；Åkerman 70\,°C 的 PPK-FIP RMSE 为 \LotoAkSeventyNew，高于 ORIG 的 \LotoAkSeventyOrig{} 和 ORIG-full 的 \LotoAkSeventyFull。

无指纹先验的联合拟合在密集采样案例中与 PPK-FIP 接近，但在两条 Åkerman 曲线中的误差更高，分别为 2.34 对 1.89 和 3.88 对 3.51。以采样时间折为单位的比较中，PPK-FIP 在 \FoldN{} 折中的 \FoldWinsOrig{} 折优于 ORIG（$p=\FoldPOrig$），在 \FoldWinsFull{} 折优于 ORIG-full（$p=\FoldPFull$）。折级结果与合并误差的改善方向一致，但不同案例与不同采样密度之间仍存在差异。
\inctab{loto_main}
\begin{figure}[!htbp]\centering\incfig{fig10_loto}\caption{留一时间点交叉验证。（a）各案例 RMSE。（b）PPK-FIP 与比较基线的 RMSE 比值。星号表示 Holm 校正后的 Wilcoxon 检验，*、**、*** 分别对应 $p<0.05$、$p<0.01$、$p<0.001$。颜色按配对误差中位差方向标注。}\label{fig:loto}\end{figure}

前半段训练的向前外推没有呈现同等幅度的合并改善（表~\ref{tab:fcst}，图~\ref{fig:fcstbars}）。PPK-FIP 在 Åkerman 100\,°C 下将 RMSE 从 \FcstAkHundredOrig{} 降至 \FcstAkHundredNew，Tao $a_w=0.25$ 下从 \FcstTaoAOrig{} 降至 \FcstTaoANew；Åkerman 70\,°C 与 Mao 的误差则高于 ORIG。合并 RMSE 为 \FcstAllNew{} 对 \FcstAllOrig{}（$p=\FcstAllPOrig$）。因此，留一时间点预测改善不能直接推广为所有训练窗口下的外推优势。

图~\ref{fig:fcst} 给出三条代表性曲线。Åkerman 100\,°C 的早期数据已包含部分减速信息，PPK-FIP 改善了后期 T 的预测；Tao $a_w=0.25$ 的早期数据不足以约束后期减速，两种机理模型均提前预测接近完全转化。Bornadel 数据中的神经 ODE 在 15\,h 后预测 M 再次上升，与观测及机理模型轨迹不同。缺少显式反应结构约束的神经 ODE 在观测窗口外产生了较大轨迹偏差。
\inctab{fcst_main}
\begin{figure}[!htbp]\centering\incfig{fig11_forecast_bars}\caption{使用 $t\le0.5\,t_\mathrm{max}$ 观测训练后的向前外推误差。图形布局与图~\ref{fig:loto} 一致。}\label{fig:fcstbars}\end{figure}
\begin{figure}[!htbp]\centering\incfig{fig12_forecast}\caption{三条反应曲线的向前外推。阴影表示训练窗口，虚线、实线和点线分别连接初始两步法、PPK-FIP 与神经 ODE 在留出时刻的预测。}\label{fig:fcst}\end{figure}

\subsection{伪持续性与指纹先验的预测贡献}
组件消融显示，伪持续性是合并留出预测改善的主要来源（表~\ref{tab:ablation}，图~\ref{fig:ablation}）。去除伪持续性后，合并 RMSE 增至 \AblNoProc{}（$p$\,\AblNoProcP），Bornadel 的误差由 \LotoBornNew{} 增至 \AblNoProcBorn，比值为 1.63。去除批内失活后，合并 LOTO RMSE 为 \AblNoKd，与完整 PPK-FIP 接近；合并外推 RMSE 为 11.41，略低于完整模型的 11.87。失活项对 Åkerman 100\,°C 的全数据拟合有帮助，但并未带来所有案例的合并预测增益。

指纹先验的作用随采样密度变化。去除指纹先验的联合估计获得合并 RMSE \AblJoint，预测损失主要出现在稀疏的 Åkerman 曲线。较紧的选择性先验 $\lambda=1$ 获得 \AblLamOne{}（$p<0.001$），主要增加 Bornadel 案例误差，说明过强约束限制了表观选择性的修正。留一案例元验证在六个案例中均选择 $\lambda=3$。在 $\lambda=3$ 下，固定 $s_p=0.05$ 或 $0.2$ 的合并 RMSE 与元选择结果一致到小数点后三位。

模型平均的贡献小于组成路径和先验的贡献。仅使用完整集总常数形式的合并 RMSE 为 \AblFullOnly，与模型平均接近；仅使用简化形式时为 \AblSimplOnly，Tao $a_w=0.25$ 的误差比值为 1.32。当前数据支持保留完整集总常数形式，但不足以说明 BIC 模型平均在其他数据条件下必然无效。
\inctab{ablation}
\begin{figure}[!htbp]\centering\incfig{fig14_ablation}\caption{PPK-FIP 的组件消融。（a）各变体合并 LOTO RMSE，虚线表示完整 PPK-FIP。（b）各案例相对于完整 PPK-FIP 的 RMSE 比值，红色表示误差增加，蓝色表示误差降低。}\label{fig:ablation}\end{figure}

跨条件参数池化未降低合并预测误差（表~\ref{tab:pooling}，图~\ref{fig:pooling}）。完全池化改善 Åkerman 100\,°C 和两个 Tao 案例，但使 Åkerman 70\,°C 的 RMSE 从 1.89 增至 2.43。留一案例元验证选择池化层级后，合并 RMSE 为 \PoolAllSel，不池化时为 \PoolAllLo。跨条件共享参数的收益取决于条件变化是否保留可共享的动力学结构，因此后续分析采用逐案例估计。
\inctab{pooling}
\begin{figure}[!htbp]\centering\incfig[0.85\linewidth]{fig17_pooling}\caption{同一反应体系不同条件之间的参数池化结果，以 LOTO RMSE 表示。}\label{fig:pooling}\end{figure}

\subsection{伪持续性路径的跨条件适用性与选择性偏差}
跨条件迁移用于检验组成路径是否具有超出单条拟合曲线的适用性。伪持续性路径在 Tao 三个水活度下的迁移误差均低于无伪持续性路径，降幅为 11--31\,\%（表~\ref{tab:transfer}，图~\ref{fig:transfer}）。Åkerman 70\,°C 与 Mao 批次的两种路径误差接近，Åkerman 100\,°C 的伪持续性路径误差较高；总体 RMSE 比值为 \TransferRatio。使用 $a_w=0.25$ 和 $0.35$ 曲线预测 $a_w=0.45$ 时，伪持续性路径减少了低反应度区域的 M 高估和 D 低估。迁移结果支持 Tao 数据中组成路径修正的跨水活度适用性，但不能单独证明酶活性位点上的微观反应过程。
\inctab{transfer}
\begin{figure}[!htbp]\centering\incfig{fig16_transfer}\caption{组成路径的跨条件迁移。（a）留出条件的组分预测误差。（b）由 $a_w=0.25$ 和 $0.35$ 预测 $a_w=0.45$ 的组成路径；虚线为无伪持续性路径，实线为伪持续性路径。}\label{fig:transfer}\end{figure}

蒙特卡洛实验量化了已知机理下的参数偏差（表~\ref{tab:sim}，图~\ref{fig:sim}）。无伪持续性真值 S1 下，PPK-FIP 与 ORIG 的曲线 RMSE 分别为 \SimSiNew{} 和 \SimSiOrig{} 摩尔\,\%，差异小于重复实验的标准差。PPK-FIP 的 $S_M$ 平均对数偏差为 $+0.025$，无正则化联合拟合的参数离散程度更大。模拟结果表明，增加机理参数在当前先验设置下未明显损害 S1 的曲线恢复精度；曲线误差接近并不等于每次拟合均将伪持续性参数置零。

伪持续性真值 S2 下，ORIG 的 $S_M$ 平均对数偏差为 \SimSMBiasOrig，对应约 \SimSMBiasOrigPct\,\% 的高估，曲线 RMSE 为 \SimSiiOrig。紧先验保留了约 $+0.22$ 的对数偏差，PPK-FIP 将偏差降至 \SimSMBiasNew、曲线误差降至 \SimSiiNew。失活真值 S3 下，PPK-FIP 的曲线误差为 \SimSiiiNew，低于 ORIG 的 \SimSiiiOrig{} 与 ORIG-full 的 \SimSiiiFull；PPK-FIP 的 $S_D$ 对数偏差仍为 $-0.129$，不能表述为完全无偏。模拟支持显式机理与适度正则化降低估计偏差，结论限于已测试的真值、噪声和采样设计。
\inctab{sim}
\begin{figure}[!htbp]\centering\incfig{fig15_simulation}\caption{三种已知真值下的蒙特卡洛估计结果。上排为相对无噪声轨迹的曲线 RMSE，下排为 $S_M$ 的对数偏差。每种情景包含 \NSimReps{} 次重复，PPK-FIP 使用 $\lambda=3$、$s_p=0.2$，紧先验 $\lambda=1$ 作为对照。}\label{fig:sim}\end{figure}

\subsection{机理约束与高斯过程残差校正的预测作用}
数据驱动基线的内插表现随采样密度变化（表~\ref{tab:lotobb}，图~\ref{fig:blackbox}a）。在稀疏的 Åkerman 曲线上，高斯过程、MLP 集成和 TabICL 的 RMSE 约为 13--21 摩尔\,\%，神经 ODE 约为 3--4，机理模型约为 1.9--3.5。Bornadel、Tao 和 Mao 的密集曲线更适合数据驱动内插。神经 ODE 的合并 LOTO RMSE 为 \LotoBBNODE，低于 PPK-FIP 的 \LotoAllNew，但合并值受各案例观测数影响，不能替代逐案例比较。

在前半段训练的外推中，四种数据驱动基线的逐案例误差均高于 PPK-FIP（表~\ref{tab:fcstbb}，图~\ref{fig:blackbox}b）。高斯过程、MLP、TabICL 和神经 ODE 的合并 RMSE 分别为 \FcstBBGP、\FcstBBMLP、\FcstBBTabICL{} 和 \FcstBBNODE，PPK-FIP 为 \FcstAllNew。本文基线未显式编码 TMP 守恒、逐级酯化顺序与 Ping Pong 速率结构，外推结果支持在当前小样本条件下保留机理约束，不能推广为所有机器学习动力学方法均缺乏外推能力。
\inctab{loto_blackbox}\inctab{fcst_blackbox}
\begin{figure}[!htbp]\centering\incfig{fig13_blackbox}\caption{机理模型与高斯过程、MLP 集成、TabICL 和神经 ODE 的预测比较，纵轴为对数尺度。（a）留一时间点交叉验证。（b）前半段训练的向前外推。}\label{fig:blackbox}\end{figure}

PPK-FIP-GP 获得了所比较方法中最低的合并 LOTO RMSE（表~\ref{tab:hybrid}，图~\ref{fig:hybrid}）。合并值为 \HybLotoAll{} 摩尔\,\%，较 ORIG 降低 \HybLotoRedOrig\,\%（$p$\,\HybLotoPOrig），在 \FoldN{} 折中的 \FoldHybWinsOrig{} 折具有较低误差。改善主要出现在密集曲线，例如 Bornadel 从 1.05 降至 0.73，Mao 从 3.10 降至 \HybLotoMao。ORIG+GP 的合并误差也降至 \HybLotoOrigGP，与 PPK-FIP-GP 的差异对应 $p=\HybLotoPOrigGP$。残差校正可改善不同机理均值的内插，因而较低内插误差本身不能唯一确认机理结构。

PPK-FIP-GP 的合并外推误差为 \HybFcstAll，与 PPK-FIP 的 \HybFcstBase{} 接近。远离训练窗口时，零均值高斯过程校正衰减，预测主要由机理均值决定。高斯过程项可吸收训练范围内未被机理模型解释的系统偏差，但偏差可能同时包含测量误差与机理缺项。本文结果支持按预测目的区分参数解释、窗口内内插与窗口外外推。
\inctab{hybrid}
\begin{figure}[!htbp]\centering\incfig{fig18_hybrid}\caption{机理模型、高斯过程残差校正模型与独立高斯过程的预测表现。（a）留一时间点交叉验证。（b）向前外推。纵轴采用对数尺度。}\label{fig:hybrid}\end{figure}

\section{讨论}
\subsection{选择性约束与动力学预测的关系}
初始两步法通过反应度域的指纹估计减少动态拟合自由度，为稀疏酯化曲线提供了参数约束 \cite{nishimura2026}。本文的组分残差、消融与模拟结果共同表明，指纹信息的作用取决于组成路径模型是否充分。Bornadel 数据中，固定不含伪持续性的指纹选择性会保留 M 与 D 分布的系统偏差；稀疏 Åkerman 数据中，完全取消指纹约束则增加预测误差。PPK-FIP 用有限宽度先验兼顾选择性稳定性与路径修正，主要贡献是改善已测试案例的留出时间点预测。

指纹先验由同一训练集估计，因此 PPK-FIP 更适合解释为数据自适应的正则化估计，而非利用独立先验信息的完整贝叶斯推断。留一时间点验证评估了估计流程的预测表现，但尚未检验后验区间的覆盖率。BIC 权重也采用近似模型权重，不应直接解释为机理真实概率。近期混合神经动力学研究同样强调参数约束和可辨识性分析的重要性 \cite{giampiccolo2024,derooij2025}；本文将参数约束具体落实为 TMP 酯化指纹估计与显式反应路径。

\subsection{伪持续性与失活参数的解释范围}
Bornadel 数据中，伪持续性修正同时改善反应度域组成路径、时间域残差和留出预测，与 Mitchell 等对传质受限条件下表观选择性的解释一致 \cite{mitchell2023}。已知真值模拟进一步显示，省略伪持续性可造成选择性偏差，适度放宽指纹先验能够减小偏差。跨水活度迁移则支持 Tao 数据中组成路径修正的可迁移性。浓度轨迹仍属于间接证据，无法唯一排除其他未建模反应或传质过程。

伪持续性概率应视为给定酶、载体与传质条件下的有效参数。Sengupta 等在另一固定化脂肪酶酯化体系中揭示了吸附和孔结构对反应过程的影响 \cite{sengupta2026}，进一步说明载体条件变化可能影响表观动力学。将 Bornadel 估计的选择性或伪持续性概率用于不同搅拌速度、颗粒尺度或载体时，需要重新验证组成路径。

批内失活主要改善 Åkerman 100\,°C 的全数据拟合，但失活项未降低合并交叉验证误差，去除失活还略微降低了合并外推误差。浓度曲线对失活的支持取决于采样窗口是否覆盖速率衰减阶段。独立酶活测量、延长时间采样及重复温度条件可帮助区分酶失活与其他减速来源。Ribeiro 等的固定化脂肪酶热失活实验支持在合适体系中考虑一阶衰减 \cite{ribeiro2025}，不能代替当前 TMP 体系的失活验证。

\subsection{参数辨识与验证设计}
初始模型核对显示，部分文献报告值与统一实现之间存在差异，多起点优化也改变了部分模型的目标值。完整集总常数模型的平坦剖面说明，在当前数据范围内，单个速率参数可能缺乏独立约束。低拟合误差、稳定参数解释和准确预测是不同评价目标，模型选择需要分别检查优化敏感性、参数区间与留出误差。

本文的留一时间点验证主要检验已观测条件和时间范围附近的预测能力。前半段训练的外推误差明显增加，且 PPK-FIP 相对于 ORIG 的合并优势不再显著，说明训练窗口的信息不足会限制机理扩展的收益。van Sluijs 等利用最优实验设计增加酶反应网络的参数信息量 \cite{vansluijs2024}，为后续 TMP 实验提供了可借鉴的方向。针对选择性与失活的实验设计，可分别增加早期 M--D 转化区间和后期减速区间的观测。

\subsection{机理模型与数据驱动校正的分工}
数据驱动方法在密集曲线内插中表现良好，说明局部预测并不总需要完整动力学参数解释。PPK-FIP-GP 在当前比较中获得最低合并内插误差，但 ORIG+GP 的表现接近，表明残差校正可以补偿不同机理模型在训练窗口内的偏差。高斯过程偏差项不能自动区分观测误差与模型缺项，因而残差校正后的高精度不构成机理正确性的独立证据。

前半段训练的外推中，PPK-FIP 优于本文设置的四种数据驱动基线。结论适用于有限样本、当前输入变量及相应训练方案。TabPFN、TabICL 和混合神经微分方程的近期研究表明，数据驱动方法的归纳约束与训练策略仍在发展 \cite{hollmann2025,qu2025tabicl,giampiccolo2024}。更明确的化学计量约束或反应速率参数化可能改变神经模型的外推表现，本文比较未覆盖所有物理约束学习方法。

\subsection{过程应用与研究局限}
PPK-FIP 可用于比较 TMP 酯化条件下的组成路径和动力学预测，但参数迁移需要保留酶--底物体系与反应条件的可比性。Tao 的跨水活度结果支持共享部分组成路径，Åkerman 的池化结果则提示温度变化可能限制参数共享。生物润滑剂过程设计应同时报告选择性估计、验证窗口和预测误差，并在新反应器条件下重新检验传质相关参数。

数据来源、验证规模和模型假设限制了结论范围。全部曲线来自已发表研究的数字化数据，缺少原始重复实验与独立误差估计，尤其无法确定 Mao 第一批次后期平台偏差的来源。四项研究形成六个拟合案例，留一案例超参数选择可能仍使用同一研究的其他条件，不能替代完全独立研究上的外部验证。边界参数的线性化置信区间也可能偏离名义覆盖率。完整参数集的实际可辨识性、先验宽度的跨体系适用性以及混合模型的守恒性质均需进一步验证。现有证据支持 PPK-FIP 在给定公开曲线上的预测改善，并支持将伪持续性和批内失活作为可检验的有效机理。

\section{结论}
本文提出 PPK-FIP，将 TMP 酯化的指纹选择性估计转化为有限宽度先验，并在共享 Ping Pong 结构中显式加入伪持续性与批内一阶失活。相对于重新估计的初始简化模型和初始全模型，PPK-FIP 的合并留一时间点预测误差分别降低 \LotoAllRedOrig\,\% 和 \LotoAllRedFull\,\%。消融结果将主要预测改善归于 Bornadel 数据的伪持续性修正，以及稀疏曲线中的指纹先验约束。已知真值模拟表明，显式伪持续性能够减小固定指纹选择性造成的估计偏差。

PPK-FIP-GP 通过高斯过程残差校正进一步降低了训练窗口内的预测误差，但前半段训练的外推中，PPK-FIP 与初始简化模型的合并误差接近。动力学模型的内插精度、外推能力与参数机理解释应分别评价。对于生物润滑剂过程设计，指纹先验和显式机理扩展可作为有限数据条件下的建模选择，参数应用仍需受数据覆盖范围、传质条件和外部验证结果约束。

\section*{数据与代码可用性}
所有数据均来自文献 \cite{nishimura2026} 的补充材料。用于解析数据、核对初始模型、实现 PPK-FIP 以及生成本文全部数字、表格和图的代码随稿件一并提供（目录 \texttt{code/}）。
\bibliographystyle{unsrtnat}
\bibliography{refs}
\end{document}
